\documentclass[11pt]{article}
\usepackage[margin=1in]{geometry}
\usepackage{amsmath,amssymb,amsthm,mathtools}
\usepackage{graphicx}
\usepackage{microtype}
\usepackage[hidelinks]{hyperref}
\usepackage{xcolor}

\newtheorem{theorem}{Theorem}
\newtheorem{lemma}{Lemma}
\newcommand{\D}{\mathbb D}
\newcommand{\N}{\mathbb N}
\newcommand{\Bmass}{\mathcal B}

\title{A strict upper bound for the bidisc Bohr radius:\\
an exact asymmetric rational-inner certificate}
\author{Automated mathematical-discovery run \texttt{fa\_banach\_001}, lane 11}
\date{12 August 2026}

\begin{document}
\maketitle

\begin{center}
\fbox{\parbox{0.91\textwidth}{\textbf{Status.}
Candidate substantial partial result, likely valid, pending human review.
The packet proves a new explicit upper bound, but it does not determine the
exact value of the bidisc Bohr radius.}}
\end{center}

\section{Source question and result}

Let
\[
 f(z,w)=\sum_{m,n\geq0}c_{mn}z^mw^n\in H^\infty(\D^2),
 \qquad
 \Bmass_r(f):=\sum_{m,n\geq0}|c_{mn}|r^{m+n}.
\]
The bidisc Bohr radius is
\[
 K_2:=\sup\{r>0:\Bmass_r(f)\leq1
       \text{ whenever }\|f\|_{H^\infty(\D^2)}\leq1\}.
\]
Knese's paper \cite{Knese} proves $K_2\geq0.3006$ and, on PDF page 4
immediately after Corollary 1.2, records the lack of a nontrivial upper bound
beyond $K_2\leq K_1=1/3$:

\begin{center}
\includegraphics[width=0.97\textwidth]{figures/open_problem_crop.png}
\end{center}

A later paper of Baran, Pikul, Woerdeman, and Wojtylak \cite{BaranEtAl}
gave $K_2<0.3177$.  The result below improves that bound.

\begin{theorem}\label{thm:main}
The two-dimensional Bohr radius satisfies
\[
 \boxed{\quad 0.3006\leq K_2<\frac{313}{1000}=0.313.\quad}
\]
More specifically, there is an explicit rational inner function
$F\in H^\infty(\D^2)$ with $\|F\|_\infty=1$ and
$\Bmass_{313/1000}(F)>1$.
\end{theorem}

\section{Proof intuition}

The one-variable extremizers are disk automorphisms applied to the coordinate
function.  On the bidisc, composing an automorphism with a \emph{symmetric}
rational inner function still approaches the one-variable threshold from the
safe side.  The useful move is to break the symmetry.  Reflect the stable
linear polynomial $1-xz-(1-x)w$ to obtain a degree-$(1,1)$ rational inner
function $\psi_x$.  At $x=1/6$, the coefficients of the resolvent
$\psi_x/(1-a\psi_x)$ have enough cancellation before absolute values are
taken that their absolute sum is unusually large.  With the exact rational
choice $a=999/1000$, only the terms through total degree eight already force
the Bohr sum above one at $r=313/1000$.

Two elementary identities make the certificate rigorous.  A sum-of-squares
identity proves directly that $\psi_x$ maps the bidisc to the unit disk.  A
four-term denominator recurrence then computes every coefficient in exact
rational arithmetic; the final strict inequality is an integer comparison.

\section{The Schur function}

Fix $0<x<1$, put $y=1-x$, and define
\begin{equation}\label{eq:psi}
 p(z,w)=1-xz-yw,
 \qquad
 \widetilde p(z,w)=zw-xw-yz,
 \qquad
 \psi_x=\frac{\widetilde p}{p}.
\end{equation}

\begin{lemma}\label{lem:schur}
The denominator $p$ has no zeros in $\D^2$, and $|\psi_x(z,w)|<1$ for
$(z,w)\in\D^2$.  Thus $\psi_x$ belongs to the bidisc Schur class.  Its
radial boundary values have modulus one almost everywhere, so
$\|\psi_x\|_\infty=1$.
\end{lemma}

\begin{proof}
If $|z|,|w|<1$, then $|xz+yw|<x+y=1$, so $p(z,w)\neq0$.
A direct expansion gives
\begin{equation}\label{eq:sos}
 |p(z,w)|^2-|\widetilde p(z,w)|^2
 =x(1-|w|^2)|1-z|^2+y(1-|z|^2)|1-w|^2.
\end{equation}
The right side is strictly positive in $\D^2$, proving $|\psi_x|<1$.
When $|z|=|w|=1$, it vanishes and $|\widetilde p|=|p|$ wherever the
quotient is defined.  Radial limits therefore have modulus one almost
everywhere, and the supremum norm is one.
\end{proof}

For $0<a<1$, let
\begin{equation}\label{eq:F}
 F_{a,x}(z,w)=\frac{a-\psi_x(z,w)}{1-a\psi_x(z,w)}.
\end{equation}
The scalar map $\zeta\mapsto(a-\zeta)/(1-a\zeta)$ is a disk automorphism.
Lemma~\ref{lem:schur} therefore shows that $F_{a,x}$ is rational inner and
$\|F_{a,x}\|_\infty=1$.

\section{Exact coefficient certificate}

Take from now on
\[
 x=\frac16,\qquad y=\frac56,\qquad
 a=\frac{999}{1000},\qquad r=\frac{313}{1000}.
\]
Write
\[
 Q:=\frac{\psi_x}{1-a\psi_x}
   =\frac{\widetilde p}{p-a\widetilde p}
   =\sum_{m,n\geq0}q_{mn}z^mw^n.
\]
Since $q_{00}=0$, equation~\eqref{eq:F} becomes
\begin{equation}\label{eq:FQ}
 F_{a,x}=a-(1-a^2)Q.
\end{equation}
The denominator of $Q$ is
\begin{equation}\label{eq:D}
 D(z,w)=1+\frac{799}{1200}z-\frac{4001}{6000}w
              -\frac{999}{1000}zw.
\end{equation}
Let $n_{11}=1$, $n_{01}=-1/6$, $n_{10}=-5/6$, and let every other
$n_{mn}$ be zero.  Comparing coefficients in $DQ=\widetilde p$ yields,
with $q_{ij}=0$ when either index is negative,
\begin{equation}\label{eq:rec}
 q_{mn}=n_{mn}-\frac{799}{1200}q_{m-1,n}
                    +\frac{4001}{6000}q_{m,n-1}
                    +\frac{999}{1000}q_{m-1,n-1}.
\end{equation}
Define the finite layer norms
\[
 T_k:=\sum_{m+n=k}|q_{mn}|.
\]
Exact application of~\eqref{eq:rec} gives
\begin{center}
\small
\begin{tabular}{c|c}
$k$ & $T_k$\\ \hline
1 & $1$\\
2 & $21983/18000$\\
3 & $18642677/12000000$\\
4 & $950105595503/648000000000$\\
5 & $4751192992371613/3888000000000000$\\
6 & $39242182643309639173/23328000000000000000$\\
7 & $2191652817952769393981/1093500000000000000000$\\
8 & $1463094271925643860042526233/
       839808000000000000000000000$
\end{tabular}
\end{center}
Consequently,
\begin{align}
 S_8&:=\sum_{k=1}^{8}T_kr^k\notag\\
 &=\frac{420207945895146052725713514433991908771800854450393}
 {839808000000000000000000000000000000000000000000000}.
 \label{eq:S8}
\end{align}
The relevant threshold is $1/(1+a)=1000/1999$, and exact subtraction gives
\begin{align}
 S_8-\frac{1000}{1999}
 &=\frac{187683844396959398701315353549825634829908046335607}
 {1678776192000000000000000000000000000000000000000000000}
 >0. \label{eq:strict}
\end{align}

By~\eqref{eq:FQ}, positivity of every omitted absolute-value term, and
~\eqref{eq:strict},
\begin{align*}
 \Bmass_r(F_{a,x})
 &\geq a+(1-a^2)\sum_{1\leq m+n\leq8}|q_{mn}|r^{m+n}\\
 &=a+(1-a^2)S_8
 >a+\frac{1-a^2}{1+a}=1.
\end{align*}
Thus the defining Bohr inequality fails at $r=313/1000$.  Since the finite
subsum depends continuously on $r$ and is strictly greater than one there,
it also fails on a left neighborhood of $313/1000$.  Hence
$K_2<313/1000$, proving Theorem~\ref{thm:main}. \qed

\section{Verification}

The included script in the \texttt{code} directory uses only Python's
standard-library exact rational-arithmetic class.  It:
\begin{enumerate}
\item recomputes all $q_{mn}$ with $1\leq m+n\leq8$ from
      recurrence~\eqref{eq:rec};
\item independently checks every coefficient residual in $DQ=\widetilde p$;
\item checks all eight displayed $T_k$ values;
\item checks~\eqref{eq:S8}, the strict threshold inequality, and the final
      positive Bohr-mass margin.
\end{enumerate}
The command is
\begin{verbatim}
conda run --no-capture-output -n sandbox python code/verify_certificate.py
\end{verbatim}
The verifier prints \texttt{PASS}.  Because all inputs and operations are
rational, this computation is an exact finite proof certificate rather than
a floating-point assumption.  Identity~\eqref{eq:sos} was also checked by
direct symbolic expansion and is short enough to audit by hand.

\section{Scope, novelty check, and reviewer focus}

The exact value of $K_2$ remains open; this theorem only narrows the interval
from the best bounds found in the search to
\[
 0.3006\leq K_2<0.313.
\]
On 12 August 2026, bounded searches covered the exact source title and arXiv
id, the phrases ``bidisc Bohr radius'', ``$K_2$ Bohr radius polydisk'',
``non-trivial upper bounds on $K_2$'', and ``asymmetric rational inner Bohr
radius bidisk'', together with the citations and later papers visible from
the source.  The separate post-source improvement found was
Baran--Pikul--Woerdeman--Wojtylak \cite[Theorem 6.4]{BaranEtAl}, which proves
$K_2<0.3177$ using a different numerically optimized degree-$(1,1)$
rational inner function and an exact degree-twelve coefficient check.
No searched source states $K_2<0.313$ or uses the one-parameter family
\eqref{eq:psi}.  The novelty assessment is therefore \emph{plausible, not
certified}; a specialist should conduct a broader bibliographic review.

The most useful human checks are: (i) expand the elementary identity
\eqref{eq:sos}; (ii) verify the sign convention in recurrence~\eqref{eq:rec};
and (iii) rerun the exact certificate.  No claim depends on the numerical scan
that suggested $x=1/6$.

\begin{thebibliography}{9}
\bibitem{Knese}
G. Knese,
\emph{Three radii associated to Schur functions on the polydisk},
arXiv:2410.21693v3 (2025).

\bibitem{BaranEtAl}
R. Baran, P. Pikul, H. J. Woerdeman, and M. Wojtylak,
\emph{Contractive realization theory for the annulus and other intersections
of discs on the Riemann sphere},
arXiv:2504.03236 (2025), Theorem 6.4.
\end{thebibliography}

\end{document}
